\documentclass[11pt]{article}
\usepackage{geometry}
\geometry{verbose}
\setlength{\parskip}{6pt}
\setlength{\parindent}{0pt}
\usepackage{color}
\usepackage{amsmath, amssymb}
\usepackage[authoryear]{natbib}
\usepackage[unicode=true,
 bookmarks=false,
 breaklinks=false,pdfborder={0 0 1},backref=section,colorlinks=false]
 {hyperref}

\makeatletter
\newcommand{\R}{\mathbb{R}}
\newcommand{\E}{\mathbb{E}}
\providecommand{\tabularnewline}{\\}
\usepackage{graphicx}
\usepackage{float}
\usepackage{url}
\usepackage{array}
\usepackage{enumitem}

\title{\Huge Assignment 1}
\author{\Large Chris Conlon}
\date{\Large Fall 2026 --- due at the start of Session 3}

\makeatother

\begin{document}

\title{Assignment 1}
\maketitle
\begin{center}
\begin{tabular*}{0.9\textwidth}{@{\extracolsep{\fill}}@{\extracolsep{\fill}}l@{\extracolsep{\fill}}l@{\extracolsep{\fill}}l}
Econometrics I & $\qquad$ & Professor Chris Conlon\tabularnewline
NYU Stern &  & Email: \href{mailto:cconlon@stern.nyu.edu}{cconlon@stern.nyu.edu}\tabularnewline
\end{tabular*}
\par\end{center}

\textbf{Instructions.} This assignment covers Lectures 1--2 only. You may use
\texttt{R} or Python. Turn in (1) a PDF presenting your results and showing your work,
and (2) your code. Derivations should use only the operator rules from Lecture 1
(linearity of $\E[\cdot]$, the variance/covariance rules, the law of iterated
expectations, and the law of total variance) --- if you find yourself writing an
integral, you are working too hard.

\subsection*{1 The algebra of expectations}

Let $X$ and $Y$ be random variables and $a,b$ constants.
\begin{enumerate}[label=(\alph*)]
\item Starting from the definition of expectation for a discrete random variable,
show that $\E\left[a + bX\right] = a + b\,\E\left[X\right]$.
\item Using only linearity (and the definition $Var(X) = \E[(X-\E X)^2]$), show that
$Var\left(a + bX\right) = b^{2}\, Var\left(X\right)$.
\item Show that $Var\left(X + Y\right) = Var\left(X\right) + Var\left(Y\right) + 2\,Cov\left(X,Y\right)$.
\item \emph{(A portfolio.)} Two assets have random returns $R_{1}, R_{2}$, each with
variance $\sigma^{2}$ and correlation $\rho$. Find the variance of the equal-weighted
portfolio return $\frac{1}{2}R_{1} + \frac{1}{2}R_{2}$. What happens as $\rho \rightarrow 1$?
As $\rho \rightarrow -1$? In one sentence: why is this the entire logic of diversification?
\item Define the prediction error $u \equiv y - \E\left[y|x\right]$. Show that for any
function $g\left(x\right)$, $\E\left[g\left(x\right) u\right] = 0$.
(Hint: law of iterated expectations.) In one sentence: why does this result mean that
$\E[y|x]$ is the ``best'' predictor of $y$ given $x$?
\end{enumerate}

\subsection*{2 Conditional expectations and total variance}

Suppose $\left(y, x\right)$ are bivariate normal with
\[
\mu = \left(\begin{array}{c} 1 \\ 2 \end{array}\right) \qquad\qquad
\Sigma = \left(\begin{array}{cc} 4 & 1 \\ 1 & 1 \end{array}\right)
\]
so $Var(y) = 4$, $Var(x) = 1$, and $Cov(y,x)=1$. Recall that for the bivariate normal,
$\E\left[y|x\right]$ is exactly linear: $\E\left[y|x\right] = \alpha + \beta x$.
\begin{enumerate}[label=(\alph*)]
\item Solve for $\alpha$ and $\beta$ using only the operator rules.
(Hint: take expectations of both sides for one equation; take the covariance of both
sides with $x$ for the other.)
\item For the bivariate normal, $Var\left(y|x\right)$ does not depend on $x$. Use the law
of total variance,
$Var\left(y\right) = \E\left[Var\left(y|x\right)\right] + Var\left(\E\left[y|x\right]\right)$,
to solve for $Var\left(y|x\right)$.
\item Show that $\frac{Var\left(\E[y|x]\right)}{Var\left(y\right)} = \rho^{2}$, where
$\rho$ is the correlation between $y$ and $x$. (This is the population $R^{2}$ of the
regression of $y$ on $x$.)
\item Now simulate 10,000 draws from this distribution. Verify (a)--(c) numerically:
regress $y$ on $x$ and compare the coefficients to (a); compare the residual variance
to (b); compare the $R^{2}$ to (c). Plot the mean of $y$ within bins of $x$ against the
fitted line.
\end{enumerate}

\subsection*{3 Watching the theorems work}

Let $X_{i} \sim \text{Exponential}(1)$ (so $\mu = 1$ and $\sigma^{2} = 1$, and the
distribution is highly skewed).
\begin{enumerate}[label=(\alph*)]
\item In Lecture 1 we derived $Var\left(\bar{X}_{n}\right) = \sigma^{2}/n$ from the
operator rules alone. For each $n \in \{10, 100, 1000\}$, simulate $S = 1{,}000$ sample
means $\bar{X}_{n}^{(s)}$ and compare the variance of the simulated sample means to
$\sigma^{2}/n$. Present the comparison in a small table.
\item Plot histograms of $\bar{X}_{n}^{(s)}$ for each $n$. Which theorem are you watching,
and what is it doing to the center and spread?
\item Now plot histograms of the \emph{standardized} means
$\sqrt{n}\left(\bar{X}_{n}^{(s)} - 1\right)$ for each $n$, overlaying a
$\mathcal{N}\left(0, 1\right)$ density. Which theorem are you watching now? At what $n$
does the normal approximation start looking trustworthy for this (skewed) distribution?
\end{enumerate}

\subsection*{4 Omitted variables, with the answer key in hand}

Recall the ``silly example'' from Lecture 2: rectangular paintings with
$\ln A_{i} = \ln W_{i} + \ln H_{i} + \varepsilon_{i}$, so the true coefficients are
$\left(1, 1\right)$ by geometry.
\begin{enumerate}[label=(\alph*)]
\item Simulate 10,000 paintings with
\[
\left(\begin{array}{c} \ln W_{i} \\ \ln H_{i} \end{array}\right) \sim
\mathcal{N}\left(\left(\begin{array}{c} 1 \\ 1 \end{array}\right),
\left(\begin{array}{cc} 1 & \sigma_{wh} \\ \sigma_{wh} & 1 \end{array}\right)\right),
\qquad \varepsilon_{i} \sim \mathcal{N}\left(0, 0.01\right)
\]
for $\sigma_{wh} = 0.6$. Run the \emph{long} regression of $\ln A$ on $\ln W$ and
$\ln H$, and the \emph{short} regression of $\ln A$ on $\ln W$ only. Report both.
\item Using the omitted variables bias formula from Lecture 2, compute what the short
regression coefficient \emph{should} be, and compare to (a).
\item Repeat (a) for $\sigma_{wh} \in \{-0.6, 0, 0.3, 0.9\}$ and plot the short-regression
coefficient against $\sigma_{wh}$, with the OVB-formula prediction overlaid as a line.
When is the short regression fine?
\item Two sentences: your colleague says ``I get a precise, highly significant
coefficient of 1.6; the standard error is tiny, so I trust it.'' What did they get
wrong, and why didn't the standard error warn them?
\end{enumerate}

\end{document}
